\documentclass[letterpaper,11pt]{article}

\usepackage{latexsym}
\usepackage[empty]{fullpage}
\usepackage{titlesec}
\usepackage{marvosym}
\usepackage[usenames,dvipsnames]{color}
\usepackage{verbatim}
\usepackage{enumitem}
\usepackage[hidelinks]{hyperref}
\usepackage{fancyhdr}
\usepackage[english]{babel}
\usepackage{tabularx}
\input{glyphtounicode}

\pagestyle{fancy}
\fancyhf{} % Clear all header and footer fields
\fancyfoot{}
\renewcommand{\headrulewidth}{0pt}
\renewcommand{\footrulewidth}{0pt}

% Adjust margins
\addtolength{\oddsidemargin}{-0.5in}
\addtolength{\evensidemargin}{-0.5in}
\addtolength{\textwidth}{1in}
\addtolength{\topmargin}{-.5in}
\addtolength{\textheight}{1.0in}

\urlstyle{same}

\raggedbottom
\raggedright
\setlength{\tabcolsep}{0in}

% Sections formatting
\titleformat{\section}{
  \vspace{-4pt}\scshape\raggedright\large
}{}{0em}{}[\color{black}\titlerule \vspace{-5pt}]

% Ensure that generated PDF is machine readable/ATS parsable
\pdfgentounicode=1

%-------------------------
% Custom commands
\newcommand{\resumeItem}[2]{
  \item\small{
    \textbf{#1}{: #2 \vspace{-2pt}}
  }
}

% For headings in list
\newcommand{\resumeHeading}[4]{
    \begin{tabular*}{0.99\textwidth}[t]{l@{\extracolsep{\fill}}r}
      \textbf{#1} & #2 \\
      \textit{\small#3} & \textit{\small #4} \\
    \end{tabular*}\vspace{-5pt}
}

\newcommand{\resumeSubheading}[4]{
  \vspace{-1pt}\item
    \begin{tabular*}{0.97\textwidth}[t]{l@{\extracolsep{\fill}}r}
      \textbf{#1} & #2 \\
      \textit{\small#3} & \textit{\small #4} \\
    \end{tabular*}\vspace{-5pt}
}

\newcommand{\resumeSubSubheading}[2]{
    \begin{tabular*}{0.97\textwidth}{l@{\extracolsep{\fill}}r}
      \textit{\small#1} & \textit{\small #2} \\
    \end{tabular*}\vspace{-5pt}
}

\newcommand{\resumeSubItem}[2]{\resumeItem{#1}{#2}\vspace{-4pt}}

\renewcommand{\labelitemii}{$\circ$}

\newcommand{\resumeSubHeadingListStart}{\begin{itemize}[leftmargin=*]}
\newcommand{\resumeSubHeadingListEnd}{\end{itemize}}
\newcommand{\resumeItemListStart}{\begin{itemize}}
\newcommand{\resumeItemListEnd}{\end{itemize}\vspace{-5pt}}

%-------------------------------------------
%%%%%%  CV STARTS HERE  %%%%%%%%%%%%%%%%%%%%%%%%%%%%

\begin{document}

%----------HEADING-----------------
\begin{tabular*}{\textwidth}{l@{\extracolsep{\fill}}r}
	\textbf{\Large Giacomo Cappelletto} & \href{mailto:giacomo.cappelletto@icloud.com}{giacomo.cappelletto@icloud.com} \\
	Boston, MA \textbar{} Treviso, Italy & \href{tel:+393246814985}{+39 324 681 4985} \textbar{} \href{tel:+18577530133}{+1 (857) 753-0133}\\
	\href{https://www.linkedin.com/in/giacomo-cappelletto/}{linkedin.com/in/giacomo-cappelletto} & \href{https://github.com/JJCAPPE}{github.com/JJCAPPE} \\
	\href{https://cv-nu-sage.vercel.app/}{cv-nu-sage.vercel.app} &
\end{tabular*}

\vspace{0.1in}

%----------PROFESSIONAL SUMMARY-----------
\section*{Riepilogo Professionale}
Studente di Ingegneria Informatica alla Boston University (GPA: 3.96) e atleta della squadra Division I di canottaggio. Ho sviluppato prodotti full-stack, app desktop native e pipeline di computer vision/biomeccanica; disponibile per internship in SWE, full-stack e ML applicato.
%-----------EDUCATION-----------------
\section{Formazione}
\resumeSubHeadingListStart
\resumeSubheading
{Boston University, Boston, MA}{2024 -- 2028}
{Laurea Quadriennale (B.Sc.) in Ingegneria Informatica; GPA: 3.96}{Freshman Student Athlete of the Year 2024-25}
\resumeSubheading
{H-Farm International School, Venice, Italy}{2022 -- 2024}
{Diploma International Baccalaureate (IB)}{HL: Matematica, Fisica, Informatica}
\resumeSubHeadingListEnd

%-----------EXPERIENCE-----------------
\section{Esperienza}
\resumeSubHeadingListStart

\resumeSubheading
{Boston University College of Engineering, Boston, MA}{Dec 2025 -- Presente}
{Ricercatore Universitario (Undergraduate Researcher)}{}
\resumeItemListStart
\resumeItem{Pipeline di biomeccanica}{Sviluppo di una pipeline in Python per stabilizzare video di voga, eseguire inferenza 2D dei keypoint su singolo atleta, ricostruire scheletri 3D e produrre angoli/posizioni articolari per frame.}
\resumeItem{Analisi cinematica}{Calcolo di cinematiche mediate sul colpo da sequenze di scheletri 3D per quantificare coordinazione, range di movimento e consistenza tecnica tra colpi.}
\resumeItem{Modellazione}{Addestramento di un modello di regressione sequence-to-sequence che mappa cinematiche allineate nel tempo in curve di forza della voga, validato contro telemetria strumentata in acqua/ergometro.}
\resumeItemListEnd

\resumeSubheading
{SOCIETA CAPPELLETTO S.R.L., Treviso, Italy}{May 2025 -- Sep 2025}
{Sviluppatore Full-Stack (Ricostruzione Sistema Inventario in Tauri/React/Rust)}{}
\resumeItemListStart
\resumeItem{Migrazione \& Prestazioni}
{Ricostruita un’app di inventario in Electron usando Tauri (Rust) + React, riducendo la dimensione dell’app del 70\%. Parallelizzate le chiamate API, migliorando la ricerca del 43\%.}
\resumeItem{Funzionalità avanzate}
{Implementati storico audit (Firebase), riconciliazione in tempo reale dell’inventario multi-sede e query Shopify ottimizzate in GraphQL; aggiunta ricerca istantanea SKU con percorsi di fallback.}
\resumeItem{Impatto operativo}
{Ridotto l’utilizzo di memoria dell’80\% tramite backend nativo in Rust; aumentata l’affidabilità ed eliminati i ritardi di ricerca con aggiornamenti immediati dell’inventario.}
\resumeItemListEnd

\resumeSubheading
{TickIT GmbH, Boston, MA (Remote)}{Oct 2024 -- Jan 2025}
{Ingegnere Full-Stack (Startup Piattaforma di Ticketing)}{}
\resumeItemListStart
\resumeItem{Panoramica}
{Sviluppate UI responsive (Next.js, shadcn/ui, Tailwind) e funzionalità backend (Rails, PostgreSQL); rilasciate integrazioni API scalabili per una piattaforma di ticketing.}
\resumeItem{Contributo}
{Contribuito all’architettura per l’autenticazione dei biglietti e migliorati i flussi Agile durante il ramp-up su un nuovo stack.}
\resumeItemListEnd

\resumeSubheading
{SOCIETA CAPPELLETTO S.R.L., Treviso, Italy}{May 2024 -- Jul 2024}
{Sviluppatore Full-Stack (Sistema di Gestione Inventario)}{}
\resumeItemListStart
\resumeItem{Progetto}
{Sviluppata un’app Electron + React + Node.js integrata con Shopify API + Firebase Firestore per sincronizzare l’inventario online su due sedi.}
\resumeItem{Impatto}
{Eliminate discrepanze di inventario e ridotto dell’80\% il lavoro manuale in cassa.}
\resumeItemListEnd

\resumeSubheading
{H-Farm International School, Treviso, Italy}{Mar 2024 -- Apr 2024}
{Sviluppatore Python (App di Pianificazione Incontri)}{}
\resumeItemListStart
\resumeItem{Applicazione}
{Sviluppata un’app Python/Streamlit che legge gli orari dal database della scuola e genera automaticamente bozze di email di richiesta incontro usando LLM locali (Ollama).}
\resumeItem{Collaborazione}
{Collaborato con docenti per adattare i flussi di lavoro e integrare con i processi esistenti, migliorando l’efficienza della pianificazione.}
\resumeItemListEnd

\resumeSubheading
{DLF Treviso, Treviso, Italy}{Sep 2023 -- Nov 2023}
{Sviluppatore iOS (App Logbook per Voga)}{}
\resumeItemListStart
\resumeItem{Sviluppo}
{Sviluppata un’app iOS in SwiftUI per il training log con Swift Charts per monitoraggio sessioni e analisi delle prestazioni.}
\resumeItem{Processo}
{Gestito lo sviluppo end-to-end e iterato sulla base del feedback degli atleti.}
\resumeItemListEnd

\resumeSubheading
{E-TRASH, Treviso, Italy}{Feb 2023 -- Jun 2023}
{Stagista Progettista CAD}{}
\resumeItemListStart
\resumeItem{Modellazione CAD}
{Progettato il primo prototipo di un cestino intelligente in Autodesk Fusion 360, con focus su integrità strutturale.}
\resumeItem{Ottimizzazione}
{Applicata ottimizzazione topologica e test ai carichi; realizzati render in Blender per presentazioni e iterazione.}
\resumeItemListEnd
\resumeSubHeadingListEnd

%-----------PROJECTS-----------------
\section{Progetti}
\resumeSubHeadingListStart
\resumeSubItem{NoteWorthy – Piattaforma di Conversione Automatica Appunti}
{Sviluppato \href{https://noteworthy-git-main-giacomo-cappellettos-projects.vercel.app/}{noteworthy.vercel.app}, un’app Next.js/TypeScript che converte appunti scritti a mano in PDF formattati e LaTeX. Integrati Google AI Studio, GitHub/Google OAuth e un compilatore LaTeX containerizzato con Docker; deploy tramite Cloud Run e Vercel.}
\resumeSubItem{Riconoscimento Immagini CIFAR-10 con Neural Architecture Search}
{Sviluppato un classificatore CIFAR-10 in TensorFlow/Keras e una pipeline di neural architecture search per generare e valutare candidati CNN. Automatizzati training, logging delle metriche e selezione del best model per training esteso.}
\resumeSubHeadingListEnd

%-----------SKILLS-----------------
\section{Competenze}
\resumeSubHeadingListStart
\item{
            \textbf{Linguaggi}{: TypeScript/JavaScript, Python, Rust, Ruby, C/C++, Swift, Java} \\
            \textbf{Framework \& Librerie}{: Next.js, React, Ruby on Rails, Node.js, Tauri, SwiftUI, TensorFlow/Keras} \\
            \textbf{Strumenti \& Piattaforme}{: Git/GitHub, Tailwind CSS, PostgreSQL, Docker, Firebase, Shopify API (GraphQL), Streamlit, Ollama, Google Cloud Run, Vercel}
      }
\resumeSubHeadingListEnd

\end{document}
